Questa sezione descrive il modello in termini astratti, prima dei dati e delle regole di dettaglio trattati nelle sezioni successive. Il sistema simulato è il servizio ferroviario regionale e suburbano della Lombardia in una giornata; il modello ne considera una parte, cioè la circolazione dei treni di Trenord sull'infrastruttura, e ne esclude i passeggeri, i treni merci, quelli a lunga percorrenza e i guasti. \\
Il modello è un sistema multi-agente situato su un grafo: gli agenti sono i treni, l'ambiente è la rete con i suoi binari, e i treni interagiscono soltanto attraverso l'occupazione dei binari.

\subsection{Ambiente}
L'ambiente è un grafo orientato. I nodi sono i punti in cui un binario inizia, finisce o si dirama; i link sono tratti di binario percorribili in un solo verso. Ogni link ha una lunghezza, una velocità massima e una capienza, cioè il numero di treni che può contenere insieme. I link sono raggruppati in risorse, come descritto nella sezione precedente: una risorsa è l'insieme dei link che vengono occupati insieme, ed è l'unità con cui il modello rappresenta una sezione di blocco, un binario unico percorso nei due sensi, un binario di stazione o l'insieme dei deviatoi all'ingresso di una stazione. In ogni istante una risorsa è libera, in uso oppure esaurita, quando ha raggiunto la propria capienza.

La \cref{fig:modello-grafo} mostra questi elementi su un tratto di rete di esempio: due stazioni descritte binario per binario, una linea a doppio binario fra di esse e una tratta a binario unico. Quattro treni sono rappresentati nelle situazioni tipiche: uno fermo a un binario di stazione, uno in corsa sulla linea con il tratto prenotato davanti a sé, uno che percorre il binario unico e uno che lo attende in stazione per l'incrocio.

\begin{figure}[H]
	\centering
	% Grafo di esempio dell'ambiente del modello: due stazioni descritte binario per binario,
% una linea a doppio binario mesoscopica e una tratta a binario unico, con tre treni.
% Incluso da modello.tex dentro un ambiente figure.
\begin{tikzpicture}[x=0.95cm, y=0.95cm, >=Stealth,
	nodo/.style={circle, draw, fill=white, inner sep=1.4pt},
	link/.style={->, line width=0.7pt},
	bidi/.style={<->, line width=0.7pt},
	risorsa/.style={draw=black!55, dashed, rounded corners=3pt},
	treno/.style={line width=3.2pt, line cap=round, color=red!75!black},
	prenotato/.style={line width=3.2pt, color=red!75!black, opacity=0.3},
	etichetta/.style={font=\scriptsize, align=center}]

	% risorse
	\draw[risorsa] (1.45,0.5) rectangle (2.95,1.1);
	\draw[risorsa] (1.45,-1.1) rectangle (2.95,-0.5);
	\draw[risorsa] (9.45,0.5) rectangle (10.95,1.1);
	\draw[risorsa] (9.45,-1.1) rectangle (10.95,-0.5);
	\draw[risorsa] (12.45,-0.35) rectangle (15.75,0.35);

	% nodi
	\node[nodo] (a0) at (0,0) {};
	\node[nodo] (a1u) at (1.2,0.8) {};
	\node[nodo] (a2u) at (3.2,0.8) {};
	\node[nodo] (a1d) at (1.2,-0.8) {};
	\node[nodo] (a2d) at (3.2,-0.8) {};
	\node[nodo] (a3) at (4.4,0) {};
	\node[nodo] (b0) at (8.2,0) {};
	\node[nodo] (b1u) at (9.2,0.8) {};
	\node[nodo] (b2u) at (11.2,0.8) {};
	\node[nodo] (b1d) at (9.2,-0.8) {};
	\node[nodo] (b2d) at (11.2,-0.8) {};
	\node[nodo] (b3) at (12.2,0) {};
	\node[nodo] (c0) at (16,0) {};

	% stazione A: un verso per binario
	\draw[link] (a0) -- (a1u);
	\draw[link] (a1u) -- (a2u);
	\draw[link] (a2u) -- (a3);
	\draw[link] (a3) -- (a2d);
	\draw[link] (a2d) -- (a1d);
	\draw[link] (a1d) -- (a0);

	% linea a doppio binario, mesoscopica
	\draw[link] ([yshift=3pt]a3.east) -- node[above, etichetta] {capienza 2} ([yshift=3pt]b0.west);
	\draw[link] ([yshift=-3pt]b0.west) -- node[below, etichetta] {capienza 2} ([yshift=-3pt]a3.east);

	% stazione B: binari percorribili nei due versi
	\draw[bidi] (b0) -- (b1u);
	\draw[bidi] (b1u) -- (b2u);
	\draw[bidi] (b2u) -- (b3);
	\draw[bidi] (b0) -- (b1d);
	\draw[bidi] (b1d) -- (b2d);
	\draw[bidi] (b2d) -- (b3);

	% binario unico
	\draw[bidi] (b3) -- node[below=11pt, etichetta] {capienza 1} (c0);

	% treni
	\draw[treno] (1.7,0.8) -- (2.7,0.8);
	\begin{scope}[yshift=3pt]
		\draw[prenotato] (5.9,0) -- (7.6,0);
		\draw[treno] (4.9,0) -- (5.9,0);
	\end{scope}
	\draw[treno] (9.7,-0.8) -- (10.7,-0.8);
	\draw[prenotato] (12.6,0) -- (13.7,0);
	\draw[treno] (13.7,0) -- (14.7,0);

	% etichette
	\node[etichetta] at (2.2,2.15) {Stazione A\\(microscopica)};
	\node[etichetta] at (6.3,2.15) {Linea a doppio binario\\(mesoscopica)};
	\node[etichetta] at (10.2,2.15) {Stazione B\\(microscopica)};
	\node[etichetta] at (14.1,2.15) {Binario unico\\(un blocco)};
	\node[etichetta] at (2.2,1.35) {binario 1};
	\node[etichetta] at (2.2,-1.4) {binario 2};
	\node[etichetta] at (10.2,1.35) {binario 1};
	\node[etichetta] at (10.2,-1.4) {binario 2};

	% legenda
	\begin{scope}[yshift=-2.3cm]
		\node[nodo] at (0.3,0) {};
		\node[etichetta, anchor=west] at (0.5,0) {nodo};
		\draw[link] (2.0,0) -- (2.9,0);
		\node[etichetta, anchor=west] at (3.0,0) {link};
		\draw[risorsa] (4.4,-0.2) rectangle (5.3,0.2);
		\node[etichetta, anchor=west] at (5.4,0) {risorsa};
		\draw[treno] (7.4,0) -- (8.2,0);
		\node[etichetta, anchor=west] at (8.4,0) {treno};
		\draw[prenotato] (10.2,0) -- (11.0,0);
		\node[etichetta, anchor=west] at (11.2,0) {tratto prenotato};
	\end{scope}
\end{tikzpicture}

	\caption{L'ambiente del modello su un tratto di rete di esempio. Ogni binario di stazione e il blocco a binario unico sono una risorsa con capienza 1; i link della linea mesoscopica ammettono due treni per verso. Le frecce doppie indicano due link di verso opposto che condividono la risorsa.}
	\label{fig:modello-grafo}
\end{figure}

La \cref{fig:modello-cadorna} mostra un caso reale, il nodo formato dalle stazioni di Milano Cadorna, Milano Domodossola e Milano Bovisa, come è descritto nel modello. Cadorna è una stazione di testa con dieci binari, ciascuno una risorsa a sé, riuniti in cinque gruppi secondo le linee che li usano. I binari non sono collegati tutti allo stesso modo: quelli da 1 a 5 raggiungono soltanto il fascio di binari del ramo di Saronno, il 6 e il 7 soltanto quello del ramo di Asso, e i binari da 8 a 10 entrambi. I deviatoi che collegano i binari a ciascun fascio formano una gola, lunga 650 metri e percorsa a 30 chilometri all'ora, che nel modello è una sola risorsa: un treno che la impegna blocca gli altri treni dello stesso fascio. Oltre la gola ogni fascio ha un binario per verso, e a Domodossola i quattro binari sono passanti, uno per fascio e per verso.

A Bovisa i due fasci provenienti da Cadorna incontrano le linee del Passante. Gli otto binari sono riuniti in quattro gruppi di due, uno per verso: i binari 1 e 2 ricevono i treni del Passante e quelli di Porta Garibaldi diretti al ramo di Saronno, il 3 e il 4 i treni di Cadorna dello stesso ramo, il 5 e il 6 i treni di Cadorna del ramo di Asso, il 7 e l'8 i treni del Passante dello stesso ramo. I binari 5 e 6 possono fare da riserva per il ramo di Saronno, e sono per questo collegati a entrambi i fasci e a entrambe le uscite. I deviatoi a sud e a nord della stazione sono due gole, ciascuna una sola risorsa comune a tutti i gruppi: è qui che i treni di Cadorna e quelli del Passante, pur avendo binari distinti, si contendono il passaggio.

\begin{figure}[H]
	\centering
	% Grafo del nodo Milano Cadorna - Milano Domodossola - Milano Bovisa come e' descritto nel modello
% (data/nodes/cadorna-bovisa.json): dieci binari di testa a Cadorna in cinque gruppi, due gole,
% due fasci di due binari, quattro binari passanti a Domodossola, otto a Bovisa in quattro gruppi,
% con gli ingressi dal Passante e da Porta Garibaldi e le uscite verso Saronno e verso Asso.
% Incluso da modello.tex dentro un ambiente figure.
\begin{tikzpicture}[x=0.8cm, y=0.8cm, >=Stealth,
	nodo/.style={circle, draw, fill=white, inner sep=1.2pt},
	snodo/.style={circle, draw, fill=black!70, inner sep=1.8pt},
	link/.style={->, line width=0.5pt},
	bidi/.style={<->, line width=0.5pt},
	riserva/.style={->, line width=0.5pt, dashed, color=black!55},
	stazione/.style={draw=black!45, line width=0.4pt, rounded corners=3pt},
	esterno/.style={rectangle, draw=black!60, fill=white, rounded corners=2pt, inner sep=2.5pt, font=\scriptsize, align=center},
	binario/.style={<->, line width=1.5pt},
	passante/.style={->, line width=1.5pt},
	etichetta/.style={font=\scriptsize, align=center},
	nota/.style={font=\scriptsize, align=left, anchor=west},
	mxp/.style={color=red!70!black},
	saronno/.style={color=blue!65!black},
	asso/.style={color=green!45!black},
	misti/.style={color=orange!85!black},
	essetre/.style={color=violet!80!black},
	tunnel/.style={color=cyan!60!black}]

	% gole: ogni area ombreggiata e' una sola risorsa
	\fill[blue!60!black, opacity=0.10] (1.8,4.7) -- (1.8,-2.45) -- (4.2,3.0) -- cycle;
	\fill[green!45!black, opacity=0.12] (1.8,0.95) -- (1.8,-2.45) -- (4.2,-0.75) -- cycle;
	\fill[black, opacity=0.07] (9.9,-3.05) rectangle (11.2,5.4);
	\fill[black, opacity=0.07] (13.4,-3.05) rectangle (14.6,5.4);

	% riquadri delle tre stazioni
	\draw[stazione] (-0.75,-2.65) rectangle (2.1,4.9);
	\draw[stazione] (6.45,-1.75) rectangle (8.65,4.0);
	\draw[stazione] (11.2,-3.05) rectangle (13.4,5.4);

	% Cadorna: binario k a y = 4.5 - 0.75 (k - 1)
	\foreach \k/\stile in {1/mxp, 2/saronno, 3/saronno, 4/saronno, 5/saronno,
			6/asso, 7/asso, 8/misti, 9/misti, 10/essetre} {
		\pgfmathsetmacro{\y}{4.5 - 0.75 * (\k - 1)}
		\node[nodo] (t\k) at (0,\y) {};
		\node[nodo] (p\k) at (1.8,\y) {};
		\draw[binario, \stile] (t\k) -- (p\k);
		\node[etichetta, anchor=east] at (-0.15,\y) {\k};
	}

	% gole di Cadorna
	\node[snodo] (g1) at (4.2,3.0) {};
	\node[snodo] (g2) at (4.2,-0.75) {};
	\foreach \k in {1,...,5} { \draw[bidi] (p\k) -- (g1); }
	\foreach \k in {6,7} { \draw[bidi] (p\k) -- (g2); }
	\foreach \k in {8,9,10} { \draw[bidi] (p\k) -- (g1); \draw[bidi] (p\k) -- (g2); }

	% Domodossola: quattro binari passanti, uno per verso
	\foreach \k/\y/\stile in {1/3.4/saronno, 2/2.6/saronno, 3/-0.35/asso, 4/-1.15/asso} {
		\node[nodo] (d\k a) at (6.8,\y) {};
		\node[nodo] (d\k b) at (8.3,\y) {};
	}
	\draw[passante, saronno] (d1a) -- (d1b);
	\draw[passante, saronno] (d2b) -- (d2a);
	\draw[passante, asso] (d3a) -- (d3b);
	\draw[passante, asso] (d4b) -- (d4a);
	\node[etichetta] at (7.55,3.67) {1};
	\node[etichetta] at (7.55,2.33) {2};
	\node[etichetta] at (7.55,-0.08) {3};
	\node[etichetta] at (7.55,-1.42) {4};

	% fasci fra Cadorna e Domodossola
	\draw[link] (g1) -- (d1a);
	\draw[link] (d2a) -- (g1);
	\draw[link] (g2) -- (d3a);
	\draw[link] (d4a) -- (g2);

	% Bovisa: otto binari passanti, dispari verso nord e pari verso Milano
	\foreach \k/\y/\stile in {1/4.9/tunnel, 2/4.2/tunnel, 3/3.4/saronno, 4/2.6/saronno,
			5/-0.35/asso, 6/-1.15/asso, 7/-2.0/tunnel, 8/-2.7/tunnel} {
		\node[nodo] (b\k a) at (11.4,\y) {};
		\node[nodo] (b\k b) at (13.2,\y) {};
	}
	\foreach \k/\stile in {1/tunnel, 3/saronno, 5/asso, 7/tunnel} {
		\draw[passante, \stile] (b\k a) -- (b\k b);
	}
	\foreach \k/\stile in {2/tunnel, 4/saronno, 6/asso, 8/tunnel} {
		\draw[passante, \stile] (b\k b) -- (b\k a);
	}
	\foreach \k/\y in {1/5.15, 3/3.65, 5/-0.1, 7/-1.75} {
		\node[etichetta] at (12.3,\y) {\k};
	}
	\foreach \k/\y in {2/3.95, 4/2.35, 6/-1.4, 8/-2.95} {
		\node[etichetta] at (12.3,\y) {\k};
	}

	% fasci fra Domodossola e Bovisa
	\draw[link] (d1b) -- (b3a);
	\draw[link] (b4a) -- (d2b);
	\draw[link] (d3b) -- (b5a);
	\draw[link] (b6a) -- (d4b);
	\draw[riserva] (d1b) -- (b5a);
	\draw[riserva] (b6a) -- (d2b);

	% ingressi da sud: Passante e Porta Garibaldi
	\node[esterno] (s1) at (7.55,5.95) {Lancetti (Passante)\\e Porta Garibaldi};
	\node[esterno] (s2) at (7.55,-3.3) {Lancetti (Passante)};
	\draw[link] (s1) -- (b1a);
	\draw[link] (b2a) -- (s1);
	\draw[link] (s2) -- (b7a);
	\draw[link] (b8a) -- (s2);

	% uscite a nord: ramo di Saronno e ramo di Asso
	\node[snodo] (ns) at (14.6,3.75) {};
	\node[snodo] (na) at (14.6,-1.55) {};
	\foreach \k in {1,3} { \draw[link] (b\k b) -- (ns); }
	\foreach \k in {2,4} { \draw[link] (ns) -- (b\k b); }
	% i binari 5 e 6 sono del ramo di Asso: verso Saronno sono la riserva
	\draw[riserva] (b5b) to[bend right=14] (ns);
	\draw[riserva] (ns) to[bend left=30] (b6b);
	\foreach \k in {5,7} { \draw[link] (b\k b) -- (na); }
	\foreach \k in {6,8} { \draw[link] (na) -- (b\k b); }
	\draw[bidi] (ns) -- (15.5,3.75);
	\draw[bidi] (na) -- (15.5,-1.55);
	\node[nota] at (15.55,3.75) {Quarto Oggiaro,\\Bollate, Saronno};
	\node[nota] at (15.55,-1.55) {Affori\\(ramo di Asso)};

	% titoli
	\node[etichetta] at (0.7,5.45) {\textbf{Milano Cadorna}\\10 binari di testa};
	\node[etichetta] at (7.55,4.5) {\textbf{Milano Domodossola}\\4 binari passanti};
	\node[etichetta] at (12.3,5.95) {\textbf{Milano Bovisa}\\8 binari passanti};
	\node[etichetta] at (10.55,-3.35) {gola sud};
	\node[etichetta] at (14.0,-3.35) {gola nord};
	\node[etichetta] at (3.4,4.3) {gola f1\\650 m, 30 km/h};
	\node[etichetta] at (3.9,-2.35) {gola f2\\650 m, 30 km/h};
	\node[etichetta] at (5.6,2.0) {fascio f1\\1,6 km};
	\node[etichetta] at (5.9,-1.75) {fascio f2\\1,6 km};
	\node[etichetta] at (9.1,3.65) {2,4 km};

	% legenda
	\begin{scope}[yshift=-3.75cm]
		\draw[binario, mxp] (0,0) -- (0.9,0);
		\node[nota] at (1.0,0) {Cadorna 1: Malpensa Express};
		\draw[binario, saronno] (6.6,0) -- (7.5,0);
		\node[nota] at (7.6,0) {ramo di Saronno};
		\draw[binario, asso] (12.0,0) -- (12.9,0);
		\node[nota] at (13.0,0) {ramo di Asso};
		\draw[binario, misti] (0,-0.6) -- (0.9,-0.6);
		\node[nota] at (1.0,-0.6) {Cadorna 8--9: variabili};
		\draw[binario, essetre] (6.6,-0.6) -- (7.5,-0.6);
		\node[nota] at (7.6,-0.6) {Cadorna 10: S3};
		\draw[binario, tunnel] (12.0,-0.6) -- (12.9,-0.6);
		\node[nota] at (13.0,-0.6) {Passante};
		\fill[black, opacity=0.12] (0,-1.4) rectangle (0.9,-1.0);
		\node[nota] at (1.0,-1.2) {gola: una sola risorsa};
		\draw[riserva] (6.6,-1.2) -- (7.5,-1.2);
		\node[nota] at (7.6,-1.2) {collegamento di riserva};
	\end{scope}
\end{tikzpicture}

	\caption{Il nodo Milano Cadorna, Milano Domodossola, Milano Bovisa nel modello. Ogni binario di stazione è una risorsa con capienza 1; le aree ombreggiate sono le gole, ciascuna una sola risorsa. I punti neri riuniscono i link di una gola o di una diramazione. Ogni binario di Cadorna è inoltre collegato al ricovero, non disegnato.}
	\label{fig:modello-cadorna}
\end{figure}

La rete è descritta a due scale, riassunte nella \cref{tab:modello-reti}. Le linee fra una stazione e l'altra sono mesoscopiche: un link rappresenta tutti i binari della tratta e la sua capienza ne esprime il numero. Le stazioni principali, i bivi e le tratte a binario unico sono microscopiche: ogni link è un singolo binario con capienza 1. Il modello è quindi ibrido, ma la differenza fra le due scale riguarda solo l'ambiente: il treno è in entrambe un individuo, con posizione e velocità proprie.

\begin{table}[H]
	\centering
	\caption{Le reti del modello. Il motore esegue la rete di dettaglio, a cui ogni simulazione aggiunge i link di fermata e di ricovero.}
	\label{tab:modello-reti}
	\begin{tabularx}{\textwidth}{Xrrr}
		\toprule
		Rete & Nodi & Link & Risorse \\
		\midrule
		Mesoscopica: una stazione, un nodo & 459 & 1195 & 309 \\
		Di dettaglio: stazioni principali binario per binario & 6008 & 11452 & 962 \\
		\bottomrule
	\end{tabularx}
\end{table}

L'ambiente non è uno sfondo passivo. Fornisce a ogni treno la sola informazione che gli serve, cioè se il tratto che ha davanti è disponibile; riceve le richieste di occupazione; decide chi ottiene una risorsa contesa; aggiorna lo stato delle risorse. Secondo la classificazione usuale degli ambienti, quello del modello è inaccessibile, perché un treno conosce lo stato delle sole risorse sul proprio percorso e non quello della rete; deterministico, perché nessuna grandezza è estratta a caso e lo stesso scenario produce sempre lo stesso esito; sequenziale, perché una risorsa occupata ora condiziona ciò che accade dopo; dinamico, perché lo stato delle risorse cambia per effetto degli altri treni. Il determinismo è una proprietà del simulatore e non del sistema reale, ed è discusso fra i limiti.

\subsection{Agenti}
Gli agenti sono i treni. In una giornata feriale il modello ne contiene 343, che effettuano 2349 corse su 57 linee. Ogni treno è descritto da tre gruppi di grandezze, riportati nella \cref{tab:modello-agente}.

\begin{table}[H]
	\centering
	\caption{Che cosa descrive un treno nel modello.}
	\label{tab:modello-agente}
	\begin{tabularx}{\textwidth}{lXl}
		\toprule
		Gruppo & Contenuto & Origine \\
		\midrule
		Attributi fissi & lunghezza, velocità massima, accelerazione, decelerazione, posti, trazione & tipo di treno \\
		Stato & posizione della testa e della coda, velocità, accelerazione, tratto e velocità autorizzati, prossima fermata, ritardo & motore di simulazione \\
		Piano & sequenza delle corse, percorso, fermate con il binario, orari & orario generato prima della simulazione \\
		\bottomrule
	\end{tabularx}
\end{table}

Gli agenti sono eterogenei: i tipi di treno sono dieci, descritti nella \cref{sec:materiale-rotabile}, con prestazioni e lunghezze diverse, e ogni treno ha un piano proprio. Il piano però non è scelto dall'agente. A differenza delle persone simulate da MATSim, un treno non valuta e non modifica il proprio piano, e la simulazione consiste di una sola esecuzione, senza il ciclo di ripianificazione. L'unica scelta che avviene durante la corsa è la deviazione su un altro binario della stessa stazione, quando quello previsto è occupato.

Rispetto alle proprietà con cui si definisce di norma un agente, il treno è situato, perché ha una posizione nell'ambiente e ciò che può fare dipende dallo stato delle risorse che ha davanti, ed è reattivo, perché accelera, rallenta, si ferma e riparte in risposta a quello stato. È autonomo solo in parte: regola da sé il proprio moto, ma non decide percorso, fermate e orari. Non è proattivo, perché non ha scopi oltre l'esecuzione dell'orario, e non è sociale in senso stretto, perché non comunica con gli altri treni. Si tratta dunque di un agente reattivo situato, il cui comportamento complessivo è ottenuto attraverso l'ambiente più che attraverso la deliberazione. \\
Va precisato che in MATSim l'agente è formalmente il conducente del veicolo; identificare l'agente con il treno è una scelta di descrizione, che non cambia il comportamento simulato.

\subsection{Materiale rotabile}
\label{sec:materiale-rotabile}
Ogni agente appartiene a un tipo di treno, e il tipo ne fissa le caratteristiche. Il parco rappresentato è quello dell'operatore, ridotto a dieci tipi. Le caratteristiche di un tipo entrano nel modello in tre punti distinti, ed è per questo che sono raccolte qui: il moto simulato dipende da lunghezza, velocità massima, accelerazione e decelerazione; il calcolo dei consumi, descritto nella sezione sui costi, dipende da massa, numero di casse, posti e tipo di trazione; il calcolo dei costi dipende dalla trazione e dal numero dei treni in servizio.

\subsubsection{Parametri del moto}
La \cref{tab:rotabile-moto} riporta i parametri che il motore di simulazione usa per muovere i treni. La lunghezza determina quanti link un treno occupa e quando la sua coda libera una risorsa. La velocità massima limita la velocità insieme a quella del link. L'accelerazione e la decelerazione determinano i tempi di percorrenza fra due fermate e, attraverso lo spazio di frenata, la lunghezza del tratto che il treno prenota davanti a sé. I posti non incidono sul moto, perché il modello non ha passeggeri, ma servono al calcolo della massa.

\begin{table}[H]
	\centering
	\caption{Tipi di treno del modello: parametri del moto. I valori contrassegnati con * sono stime.}
	\label{tab:rotabile-moto}
	\small
	\begin{tabularx}{\textwidth}{Xlrrrrr}
		\toprule
		Tipo & Composizione & Lungh. & Posti & Vel. max & Accel. & Decel. \\
		 & & (m) & & (km/h) & (m/s\textsuperscript{2}) & (m/s\textsuperscript{2}) \\
		\midrule
		TSR & 4 casse, due piani & 105,0 & 436 & 140 & 1,00* & 0,50* \\
		TAF & 4 casse, due piani & 104,0 & 469 & 140 & 0,80* & 0,50* \\
		Caravaggio ETR 421 & 4 casse, due piani & 109,6 & 466 & 160 & 1,10 & 0,50* \\
		Caravaggio ETR 521 & 5 casse, due piani & 136,8 & 598 & 160 & 1,10 & 0,50* \\
		Donizetti ETR 204 & 4 casse, un piano & 84,2 & 262 & 160 & 1,00* & 0,50* \\
		ETR 245 Coradia Meridian & 5 casse, un piano & 82,2 & 230 & 160 & 1,00* & 0,50* \\
		ETR 425 Coradia Meridian & 5 casse, un piano & 82,2 & 290 & 160 & 1,00* & 0,50* \\
		FLIRT TSI (TILO) ETR 524 & 6 casse, un piano & 105,0 & 244 & 160 & 1,00* & 0,50* \\
		ATR 125 Stadler GTW & 4 casse, 2 moduli motore & 77,3 & 231 & 140 & 0,60* & 0,50* \\
		Colleoni ATR 803 & 3 casse, 1 modulo motore & 66,8 & 168 & 140 & 0,72* & 0,50* \\
		\bottomrule
	\end{tabularx}
\end{table}

Lunghezza, posti e velocità massima vengono dalle schede pubblicate dall'operatore e dai costruttori. L'accelerazione è pubblicata soltanto per i Caravaggio; per gli altri tipi è stimata, con valori più bassi per il TAF, il più anziano degli elettrotreni, e per i due treni diesel. Quella del Colleoni è ricavata dall'ATR 125 applicando il miglioramento del 20\% dichiarato dall'operatore rispetto alla flotta diesel precedente. La decelerazione è la stessa per tutti i tipi, \SI{0,5}{\metre\per\second\squared}, il valore predefinito di railsim: è una decelerazione di servizio prudente, e incide direttamente sulla distanza che separa due treni che si seguono.

\subsubsection{Parametri per i consumi}
La \cref{tab:rotabile-energia} riporta i parametri usati dal modello energetico. La massa in servizio è la massa a vuoto più quella dei viaggiatori, calcolata sui posti a sedere a 75 chilogrammi per persona e sulla quota di posti occupati nella fascia oraria della corsa; nella tabella è indicata quella a posti tutti occupati, che è il valore massimo usato dal modello. I servizi ausiliari, cioè climatizzazione, illuminazione e impianti di bordo, sono stimati in 30 kW per cassa. La potenza installata è riportata per completezza: il modello non la usa, perché il moto è determinato dall'accelerazione.

\begin{table}[H]
	\centering
	\caption{Tipi di treno del modello: parametri per il calcolo dei consumi. I valori contrassegnati con * sono stime.}
	\label{tab:rotabile-energia}
	\small
	\begin{tabularx}{\textwidth}{Xlrrrr}
		\toprule
		Tipo & Trazione & Massa a vuoto & Massa a pieno & Ausiliari & Potenza \\
		 & & (t) & carico (t) & (kW) & (kW) \\
		\midrule
		TSR & elettrica, 3 kV c.c. & 222 & 255 & 120 & 2720 \\
		TAF & elettrica, 3 kV c.c. & 213 & 248 & 120 & 3640 \\
		Caravaggio ETR 421 & elettrica, 3 kV c.c. & 261* & 296 & 120 & 3400 \\
		Caravaggio ETR 521 & elettrica, 3 kV c.c. & 325* & 370 & 150 & 3400 \\
		Donizetti ETR 204 & elettrica, 3 kV c.c. & 143,5 & 163 & 120 & 2000 \\
		ETR 245 Coradia Meridian & elettrica, 3 kV c.c. & 160* & 177 & 150 & -- \\
		ETR 425 Coradia Meridian & elettrica, 3 kV c.c. & 160 & 182 & 150 & 2052 \\
		FLIRT TSI (TILO) ETR 524 & elettrica, 3 kV c.c. e 15 kV c.a. & 177,5 & 196 & 180 & 2600 \\
		ATR 125 Stadler GTW & diesel-elettrica & 136 & 153 & 120 & 1160 \\
		Colleoni ATR 803 & diesel-elettrica con batterie & 135 & 148 & 90 & 1200 \\
		\bottomrule
	\end{tabularx}
\end{table}

Il tipo di trazione divide i treni in due gruppi, trattati in modo diverso sia nei consumi sia nei costi. Gli otto tipi elettrici prendono energia dalla linea di contatto e in frenata ne restituiscono una parte, che altri treni possono riusare. I due tipi diesel consumano gasolio e nel modello non recuperano energia in frenata; il Colleoni dispone di batterie, che il modello non rappresenta.

Tre masse sono stimate. La massa a vuoto dei Caravaggio non è pubblicata, ed è ricavata dal carico per asse di 18,5 tonnellate dichiarato per le motrici, esteso a tutti gli assi del convoglio per ottenere la massa a pieno carico, da cui è sottratto il peso dei viaggiatori seduti. È un limite superiore, perché le casse rimorchiate hanno verosimilmente un carico per asse inferiore: la stima dà 2,4 tonnellate per metro di lunghezza, contro 2,1 del TSR e 2,0 del TAF, i due treni a due piani con dato pubblicato. Un consumo calcolato con queste masse è quindi sovrastimato, non sottostimato. Per l'ETR 245 non è stato trovato un dato proprio, ed è adottata la massa dell'ETR 425, convoglio della stessa famiglia con lo stesso numero di casse e la stessa lunghezza.

\subsubsection{Recupero dell'energia in frenata}
In frenata i motori di un treno a trazione elettrica possono funzionare da generatori e trasformare in energia elettrica una parte dell'energia di moto del treno. La \cref{tab:rotabile-recupero} indica quali tipi, nel modello, recuperano energia in frenata e la restituiscono alla linea di contatto.

\begin{table}[H]
	\centering
	\caption{Recupero dell'energia in frenata nel modello, per tipo di treno.}
	\label{tab:rotabile-recupero}
	\small
	\begin{tabularx}{\textwidth}{lccX}
		\toprule
		Tipo & Recupero & Rendimento & Destinazione dell'energia recuperata \\
		\midrule
		TSR & sì & 0,80 & ausiliari del treno, poi linea di contatto \\
		TAF & sì & 0,80 & ausiliari del treno, poi linea di contatto \\
		Caravaggio ETR 421 & sì & 0,80 & ausiliari del treno, poi linea di contatto \\
		Caravaggio ETR 521 & sì & 0,80 & ausiliari del treno, poi linea di contatto \\
		Donizetti ETR 204 & sì & 0,80 & ausiliari del treno, poi linea di contatto \\
		ETR 245 Coradia Meridian & sì & 0,80 & ausiliari del treno, poi linea di contatto \\
		ETR 425 Coradia Meridian & sì & 0,80 & ausiliari del treno, poi linea di contatto \\
		FLIRT TSI (TILO) ETR 524 & sì & 0,80 & ausiliari del treno, poi linea di contatto \\
		ATR 125 Stadler GTW & no & -- & l'energia di frenata è dissipata \\
		Colleoni ATR 803 & no & -- & l'energia di frenata è dissipata \\
		\bottomrule
	\end{tabularx}
\end{table}

Nel modello tutti e otto i tipi elettrici recuperano, con lo stesso rendimento: della potenza frenante alle ruote, l'80\% torna disponibile come energia elettrica. Il modello non distingue quindi fra un treno recente e uno più anziano, come il TAF: è una semplificazione, adottata perché il rendimento di recupero dei singoli tipi non è pubblicato. Combinato con il rendimento di trazione, anch'esso di 0,80, il valore corrisponde a un recupero del 64\% dell'energia spesa per accelerare, in linea con il 60--70\% indicato in letteratura per i treni moderni.

L'energia recuperata non è però tutta riutilizzata. Sulle linee a 3 kV in corrente continua, che sono quelle della rete modellata, le sottostazioni non riprendono energia dalla linea di contatto: la corrente restituita da un treno in frenata è utile soltanto se nelle vicinanze un altro treno sta assorbendo potenza in quello stesso momento. Il modello segue quest'ordine: l'energia rigenerata alimenta anzitutto i servizi ausiliari del treno che frena; la parte rimanente è ceduta ai treni elettrici in trazione che si trovano entro 10 chilometri, a partire dal più vicino; ciò che nessuno assorbe è dissipato a bordo. Restituire energia alla rete significa quindi, nel modello come nella realtà di queste linee, cederla ad altri treni e non al gestore dell'energia: per questo l'energia recuperata non è contata come ricavo nei costi, ma soltanto come minore prelievo dalle sottostazioni.

I due tipi diesel non recuperano energia. La loro trasmissione è elettrica, e nella realtà una parte dell'energia di frenata può alimentare i servizi di bordo; il Colleoni dispone inoltre di batterie. Il modello non rappresenta né l'una né l'altra possibilità, e i treni diesel non partecipano allo scambio di energia sulla linea di contatto. È uno dei motivi per cui il loro consumo, senza correzione, risultava sovrastimato.

Nel giorno feriale reale i treni elettrici rigenerano in frenata 450 MWh. Di questi 307 sono riusati da altri treni e 97 sono dissipati; la parte restante alimenta i servizi ausiliari dei treni che frenano. Senza alcun recupero l'energia da prelevare dalle sottostazioni sarebbe di 1969 MWh anziché 1616: il recupero riduce il prelievo del 18\%. La quota riusata dipende dalla densità del traffico: è alta dove i treni sono vicini fra loro, e quasi nulla sulle linee isolate.

\subsubsection{I dieci tipi}
Il \textit{TSR}, Treno Servizio Regionale, è un elettrotreno a due piani a composizione variabile, da due a sei casse tutte motrici. Il modello ne rappresenta la composizione a quattro casse; in esercizio la più frequente è di otto, ottenuta accoppiando due convogli, e la capienza e la lunghezza reali possono quindi essere doppie di quelle del modello. È il treno più diffuso sulle linee suburbane.

Il \textit{TAF}, Treno ad Alta Frequentazione, è un elettrotreno a due piani a composizione bloccata di quattro casse, due motrici di estremità e due rimorchiate. È il più anziano fra i treni del modello, e le sue masse sono le meglio documentate, perché riportate nelle norme di circolazione del gestore della rete; da queste viene anche il peso di 75 chilogrammi per viaggiatore usato per tutti i tipi.

Il \textit{Caravaggio} è un elettrotreno a due piani ad alta capacità, fra i treni di più recente costruzione, presente nel modello in due versioni: ETR 421 a quattro casse e ETR 521 a cinque. Le due casse di estremità sono motrici e quelle intermedie rimorchiate. È il treno più capiente e più pesante del modello, e l'unico di cui sia pubblicata l'accelerazione.

Il \textit{Donizetti}, ETR 204, è un elettrotreno a un piano a media capacità, di quattro casse, in servizio sulle linee regionali a domanda minore, fra cui quelle della Valtellina e del lago di Como. È il più leggero fra gli elettrotreni del modello.

I \textit{Coradia Meridian} sono elettrotreni articolati a un piano, presenti in due tipi distinti: l'ETR 245, allestimento nato per il collegamento con l'aeroporto di Malpensa e oggi impiegato sulle linee suburbane, e l'ETR 425, in servizio sulle regionali. La versione a sei casse, ETR 526, non ha un tipo proprio ed è rappresentata dall'ETR 425.

Il \textit{FLIRT TSI}, ETR 524, è un elettrotreno a due tensioni di alimentazione, in servizio sulle relazioni fra la Lombardia e il Canton Ticino; nel modello circola sulla sola linea RE80.

L'\textit{ATR 125} è un'automotrice diesel-elettrica articolata, formata da due unità accoppiate, ciascuna con due casse e un modulo motore centrale. È in servizio sulle linee non elettrificate, ed è il tipo su cui è stato tarato il consumo di gasolio.

Il \textit{Colleoni}, ATR 803, è un treno diesel-elettrico di tre casse con un modulo motore centrale e batterie, anch'esso di recente costruzione, in servizio sulle linee non elettrificate; nel modello circola sulle linee R8, R35, R36 e R37.

\subsubsection{Consistenza della flotta e assegnazione alle linee}
La \cref{tab:rotabile-flotta} riporta il numero di convogli di cui l'operatore dispone per ogni tipo. Per i treni di nuova costruzione il dato viene dai comunicati dell'operatore; per gli altri è meno certo, e in parte deriva da rilievi dell'autore. Al 31 dicembre 2024 la flotta contava in tutto 390 treni elettrici e 66 diesel, oltre ai convogli formati da locomotiva e carrozze, che il modello non rappresenta.

\begin{table}[H]
	\centering
	\caption{Convogli disponibili per tipo.}
	\label{tab:rotabile-flotta}
	\begin{tabularx}{\textwidth}{lrX}
		\toprule
		Tipo & Convogli & Nota \\
		\midrule
		Caravaggio (ETR 421 e 521) & 136 & 123 a fine 2025, più 13 in consegna nel 2026 \\
		TSR & circa 104 & consegnati fra il 2007 e il 2018; dato da verificare \\
		Donizetti & 61 & \\
		ETR 425 ed ETR 526 & 54 & rilievo dell'autore \\
		TAF & 34 & rilievo dell'autore; l'operatore ne indica 25 rinnovati \\
		Colleoni & 30 & \\
		ATR 125 & 20 & sulle stesse linee circolano altre automotrici diesel, non nel modello \\
		ETR 245 & 14 & rilievo dell'autore \\
		FLIRT TSI (TILO) & 9 & quelli acquistati dal gruppo dell'operatore \\
		\bottomrule
	\end{tabularx}
\end{table}

L'operatore non assegna i treni per linea ma per impianto di manutenzione: le linee fanno capo a un impianto, e il treno di una corsa dipende dal turno, con una composizione che varia da un giorno all'altro. Il modello approssima questa situazione con una ripartizione per linea, riportata nella \cref{tab:rotabile-quote}: ogni linea ha un tipo unico oppure una composizione in percentuale, e i treni della linea sono distribuiti fra i tipi secondo quelle quote. Il tipo è attribuito al treno, cioè all'intero giro di corse che effettua nella giornata, e non alla singola corsa: le quote non cambiano quindi il numero dei treni, ma soltanto la loro ripartizione.

\begin{table}[H]
	\centering
	\caption{Ripartizione dei tipi di treno sulle linee.}
	\label{tab:rotabile-quote}
	\begin{tabularx}{\textwidth}{>{\raggedright\arraybackslash}p{0.42\textwidth}X}
		\toprule
		Linee & Tipi \\
		\midrule
		S1, S2, S5, S6, S8, S9, S12, S13, S19 & 60\% TSR, 20\% Caravaggio a 5 casse, 10\% ETR 245, 10\% TAF \\
		S3, S4 & 70\% Caravaggio a 5 casse, 30\% TSR \\
		S11 & Caravaggio a 5 casse \\
		RE51, RE54 (collegamenti con Malpensa) & Caravaggio a 4 casse \\
		RE8, RE13, R6, R7, R11, R12, R13, R21, R34 & Donizetti \\
		S7, S31, R3, R9, R18, RE3 & ATR 125 \\
		R8, R35, R36, R37 & Colleoni \\
		RE80 & FLIRT TSI (TILO) \\
		Altre linee regionali e RegioExpress & 60\% Caravaggio a 5 casse, 30\% TSR, 10\% ETR 425 \\
		\bottomrule
	\end{tabularx}
\end{table}

Le quote riflettono ciò che si osserva in esercizio: sulle linee regionali circolano soprattutto i Caravaggio, e il resto è coperto da TSR e Coradia Meridian; sulle suburbane prevalgono i TSR, mentre i TAF coprono le necessità residue. Sono state fissate in modo che i treni di ogni tipo richiesti dall'orario di un giorno feriale restino entro la consistenza della flotta; il confronto è nell'analisi dei risultati. Tipi, parametri e ripartizioni possono essere modificati dall'applicazione prima di ogni simulazione.

\subsubsection{Differenze dal materiale reale}
Il modello si discosta dal materiale reale in alcuni punti, che vanno tenuti presenti nel leggere i risultati. La composizione dei TSR è fissata a quattro casse, mentre in esercizio varia. I posti dei Caravaggio sono quelli della versione di serie, 466 e 598, mentre l'allestimento dell'operatore ne ha 443 e 573. I posti del Colleoni comprendono 17 strapuntini. Mancano del tutto i convogli formati da locomotiva e carrozze e le automotrici diesel più anziane, le cui corse sono attribuite agli altri tipi. Infine ogni tipo ha un solo insieme di prestazioni: il modello non distingue fra un treno a pieno carico e uno vuoto nel moto, ma soltanto nel calcolo dei consumi.

\subsection{Regole di comportamento}
Il comportamento di un treno segue poche regole locali, quelle di railsim descritte nella sezione precedente. A ogni aggiornamento il treno chiede il tratto di percorso che gli serve per arrestarsi alla velocità corrente; se lo ottiene accelera fino alla velocità consentita, altrimenti rallenta, si ferma prima del tratto occupato e ripete la richiesta dopo dieci secondi. Alle fermate attende l'orario di partenza, e non parte mai in anticipo. Le risorse sono concesse nell'ordine in cui vengono chieste, senza priorità fra i treni.

A queste regole il progetto ne aggiunge altre, perché quelle di base non bastano a far circolare l'orario reale sulla rete lombarda. Riguardano l'uso dei binari nelle stazioni, la precedenza e il consenso all'ingresso nelle tratte a binario unico, la prevenzione degli stalli, il ricovero dei treni fermi fra una corsa e l'altra e il loro ritiro a fine servizio. Sono descritte, ciascuna con il caso che l'ha resa necessaria, nelle \cref{sec:orari,sec:dettagli-rete}.

\subsection{Interazione e coordinamento}
I treni non si scambiano messaggi. Un treno modifica l'ambiente occupando e liberando risorse, e gli altri percepiscono la modifica quando chiedono le stesse risorse: l'interazione è indiretta e fondata sullo spazio. Il contesto di un agente è il tratto del proprio percorso compreso nello spazio di frenata, che cresce con la velocità. Un treno non sa quale altro treno occupa una risorsa, né per quanto tempo.

Nello schema con cui Ferber classifica le interazioni secondo scopi, risorse e capacità degli agenti, quella fra i treni è un'ostruzione: gli scopi sono compatibili, perché ogni treno deve soltanto eseguire il proprio orario, e le capacità sono sufficienti, ma le risorse non bastano per tutti nello stesso momento. Nessun treno ha interesse a ostacolarne un altro; si ostacolano perché usano gli stessi binari.

Il coordinamento non è però interamente locale. Tre componenti vedono lo stato di tutta la rete: il gestore delle risorse, che ne tiene lo stato; la disposizione dei treni, che concede o nega i tratti richiesti e sceglie le deviazioni; la prevenzione degli stalli, che prima di concedere l'ingresso in una tratta a binario unico considera i treni diretti verso la stessa tratta. Nel modello questi componenti fanno parte dell'ambiente, a cui spetta regolare gli effetti delle azioni e i conflitti, e corrispondono nella realtà al segnalamento e alla regolazione della circolazione. Il sistema non va quindi presentato come del tutto decentrato.

\subsection{Tempo}
Il tempo è discreto, con un passo di un secondo. All'interno del passo il motore elabora una coda di eventi ordinati per istante, uno per ogni azione di un treno: partenza, ingresso della testa in un link, uscita della coda, variazione di velocità, attesa di una prenotazione. L'aggiornamento è quindi asincrono, e l'ordine degli eventi stabilisce chi ottiene per primo una risorsa. Lo spazio è invece continuo: la posizione di un treno è una distanza in metri lungo un link, e un treno occupa in genere più link. \\
La giornata simulata inizia con la prima corsa e termina quando l'ultimo treno ha concluso il servizio, oltre la mezzanotte.

\subsection{Che cosa emerge}
Nessuna regola del modello parla di ritardi, di code o di stalli: sono effetti che nascono dall'interazione fra i treni. Un treno in ritardo occupa una risorsa più a lungo del previsto e ritarda chi lo segue; su un binario unico il ritardo passa anche ai treni del verso opposto, che attendono l'incrocio. È una retroazione positiva, contrastata dai margini dell'orario e dalle soste ai capolinea, che riassorbono il ritardo. Quando la prima prevale, davanti a una stazione satura si forma una coda, e in alcune configurazioni più treni attendono ciascuno una risorsa tenuta da un altro, senza possibilità di uscita. L'effetto di una modifica non è proporzionale alla sua entità: come mostrato nell'analisi dei risultati, il raddoppio delle corse di una linea può moltiplicare il ritardo medio per più di dieci, e la regola con cui si assegna un binario in una stazione può cambiare l'esito dell'intera giornata.

\subsection{Classificazione del modello}
La \cref{tab:modello-proprieta} riassume le proprietà del modello.

\begin{table}[H]
	\centering
	\caption{Proprietà del modello.}
	\label{tab:modello-proprieta}
	\begin{tabularx}{\textwidth}{lX}
		\toprule
		Proprietà & Valore \\
		\midrule
		Tipo & sistema multi-agente situato su grafo \\
		Agenti & treni, reattivi, eterogenei, con piano fissato \\
		Interazione & indiretta, attraverso l'occupazione delle risorse \\
		Scala & ibrida: linee mesoscopiche, stazioni principali microscopiche \\
		Tempo & discreto, passo di un secondo, eventi ordinati dentro il passo \\
		Spazio & continuo lungo i link \\
		Casualità & assente: modello deterministico \\
		Iterazioni & una \\
		Passeggeri & assenti \\
		\bottomrule
	\end{tabularx}
\end{table}

Il modello non è un automa cellulare, anche se ne ricorda un aspetto. In un automa cellulare lo spazio è un reticolo regolare di celle omogenee, un'entità è lo stato di una cella, il vicinato è fisso e tutte le celle si aggiornano insieme. Qui lo spazio è un grafo irregolare, i cui link differiscono per lunghezza, velocità e capienza; il treno è un'entità distinta dall'ambiente, con posizione e velocità continue, che occupa più link; ciò che un treno percepisce dipende dalla sua velocità; l'aggiornamento è asincrono; gli agenti sono eterogenei. Il punto di contatto è la sezione di blocco, che come una cella ha un insieme finito di stati e passa da libera a occupata.
